\documentclass[12pt,a4paper]{article}

% PAQUETES
\usepackage[utf8]{inputenc}
\usepackage[spanish]{babel}
\usepackage[margin=2.5cm]{geometry}
\usepackage{amsmath,amssymb,amsfonts}
\usepackage{graphicx}
\usepackage{textcomp}
\usepackage{xcolor}
\usepackage{float}
\usepackage{url}
\usepackage{booktabs}
\usepackage{array}
\usepackage{multirow}
\usepackage{hyperref}
\usepackage{listings}
\usepackage{matlab-prettifier}
\usepackage{listings}
\usepackage{xcolor}

\lstdefinestyle{matlabstyle}{
    language           = Matlab,
    basicstyle         = \ttfamily\footnotesize,
    keywordstyle       = \color{blue},
    commentstyle       = \color{gray},
    stringstyle        = \color{purple},
    numbers            = left,
    numberstyle        = \tiny\color{gray},
    stepnumber         = 1,
    numbersep          = 8pt,
    backgroundcolor    = \color{white},
    frame              = single,
    rulecolor          = \color{black!30},
    breaklines         = true,
    breakatwhitespace  = false,
    showstringspaces   = false,
    tabsize            = 2,
    captionpos         = b
}

\lstset{style=matlabstyle}

\begin{document}
\section{Ejercicio 3}

3. El rectificador de la figura con: Vs=208, f=60Hz V; R=2(a+b+c)$\Omega$; L=2(b+c)mH;
determinar: a) las funciones de tensión y corriente de salida (dibújelas), b) las tensiones dc y
eficaz y la corrientes dc y rms en la carga y en el diodo, d) la potencia, el factor de potencia
la eficiencia y los FR para tensión y corriente. Realizar la simulación y las comparaciones.

\begin{figure}[h]
    \centering
    \includegraphics[width=0.7\textwidth]{circuito3.png}
    \caption{Circuito 3: Rectificador monofásico de onda completa tipo puente}
\end{figure}

\subsection{Procedimiento}

Desarrollo matemático:

\begin{equation}
    R = 2(a+b+c) = 2(3+9+3) = 30{,}0\,\Omega
\end{equation}

\begin{equation}
    L = 2(b+c) = 2(9+3) = 24{,}0\,\text{mH} = 0{,}0240\,\text{H}
\end{equation}

\begin{equation}
    V_{in} = 208\,\text{V}, \quad f = 60{,}0\,\text{Hz}
\end{equation}

\begin{equation}
    V_m = 208\,\text{V} \times 1{,}414 = 294\,\text{V}
\end{equation}

\begin{equation}
    \omega = 2\pi f = 2 \times 3{,}142 \times 60{,}0\,\text{Hz} = 377\,\text{rad/s}
\end{equation}

\textbf{Para $n=0$, $m=0$:}

\begin{equation}
    V_{DC} = \frac{1}{\pi\,\text{rad}} \int_{0\,\text{rad}}^{\pi\,\text{rad}} 294\,\text{V}\sin(\theta)\,d\theta
\end{equation}

\begin{equation}
    V_{DC} = \frac{294\,\text{V}}{\pi\,\text{rad}}\left[-\cos(\theta)\right]_{0\,\text{rad}}^{\pi\,\text{rad}}
\end{equation}

\begin{equation}
    \boxed{V_{DC} = 187\,\text{V}}
\end{equation}

\begin{equation}
    I_{DC} = \frac{187\,\text{V}}{30{,}0\,\Omega} = \boxed{6{,}23\,\text{A}}
\end{equation}

\textbf{Para $n=1$, $m=2$, $f=120$\,Hz:}

\begin{equation}
    a_1 = \frac{2}{\pi\,\text{rad}} \int_{0\,\text{rad}}^{\pi\,\text{rad}} 294\,\text{V}\sin(\theta)\cos(2\theta)\,d\theta
\end{equation}

\begin{equation}
    a_1 = \frac{588\,\text{V}}{\pi\,\text{rad}} \int_{0\,\text{rad}}^{\pi\,\text{rad}} \sin(\theta)\cos(2\theta)\,d\theta
\end{equation}

\begin{equation}
    a_1 = -125\,\text{V}
\end{equation}

\begin{equation}
    b_1 = \frac{2}{\pi\,\text{rad}} \int_{0\,\text{rad}}^{\pi\,\text{rad}} 294\,\text{V}\sin(\theta)\sin(2\theta)\,d\theta = 0\,\text{V}
\end{equation}

\begin{equation}
    \omega_1 = 2\omega = 2 \times 377\,\text{rad/s} = 754\,\text{rad/s}
\end{equation}

\begin{equation}
    X_{L,1} = \omega_1 L = 754\,\text{rad/s} \times 0{,}024\,\text{H} = 18{,}1\,\Omega
\end{equation}

\begin{equation}
    Z_1 = 30{,}0\,\Omega + j18{,}1\,\Omega = 35{,}0\,\Omega \angle 0{,}543\,\text{rad}
\end{equation}

\begin{equation}
    I_1 = \frac{125\,\text{V}\angle 0\,\text{rad}}{35{,}0\,\Omega \angle 0{,}543\,\text{rad}} = 3{,}57\,\text{A}\angle(-0{,}543\,\text{rad})
\end{equation}

\textbf{Para $n=2$, $m=4$, $f=240$\,Hz:}

\begin{equation}
    a_2 = \frac{2}{\pi\,\text{rad}} \int_{0\,\text{rad}}^{\pi\,\text{rad}} 294\,\text{V}\sin(\theta)\cos(4\theta)\,d\theta = -25{,}0\,\text{V}
\end{equation}

\begin{equation}
    b_2 = \frac{2}{\pi\,\text{rad}} \int_{0\,\text{rad}}^{\pi\,\text{rad}} 294\,\text{V}\sin(\theta)\sin(4\theta)\,d\theta = 0\,\text{V}
\end{equation}

\begin{equation}
    \omega_2 = 4\omega = 4 \times 377\,\text{rad/s} = 1508\,\text{rad/s}
\end{equation}

\begin{equation}
    X_{L,2} = \omega_2 L = 1508\,\text{rad/s} \times 0{,}0240\,\text{H} = 36{,}2\,\Omega
\end{equation}

\begin{equation}
    Z_2 = 30{,}0\,\Omega + j36{,}2\,\Omega = 47{,}0\,\Omega \angle 0{,}879\,\text{rad}
\end{equation}

\begin{equation}
    I_2 = \frac{25{,}0\,\text{V}\angle 0\,\text{rad}}{47{,}0\,\Omega \angle 0{,}879\,\text{rad}} = 0{,}531\,\text{A}\angle(-0{,}879\,\text{rad})
\end{equation}

\textbf{Para $n=3$, $m=6$, $f=360$\,Hz:}

\begin{equation}
    a_3 = \frac{2}{\pi\,\text{rad}} \int_{0\,\text{rad}}^{\pi\,\text{rad}} 294\,\text{V}\sin(\theta)\cos(6\theta)\,d\theta = -10{,}7\,\text{V}
\end{equation}

\begin{equation}
    b_3 = \frac{2}{\pi\,\text{rad}} \int_{0\,\text{rad}}^{\pi\,\text{rad}} 294\,\text{V}\sin(\theta)\sin(6\theta)\,d\theta = 0\,\text{V}
\end{equation}

\begin{equation}
    \omega_3 = 6\omega = 6 \times 377\,\text{rad/s} = 2262\,\text{rad/s}
\end{equation}


\textbf{Funciones de tensión y corriente de salida:}

\begin{equation}
\begin{split}
    v_o(t) = 187\,\text{V} &- 125\,\text{V}\cos(754\,\text{rad/s}\cdot t) \\
    &- 25{,}0\,\text{V}\cos(1508\,\text{rad/s}\cdot t) \\
    &- 10{,}7\,\text{V}\cos(2262\,\text{rad/s}\cdot t)
\end{split}
\end{equation}

\begin{equation}
\begin{split}
    i_o(t) = 6{,}23\,\text{A} &+ 3{,}56\,\text{A}\sin(754\,\text{rad/s}\cdot t - 0{,}543\,\text{rad}) \\
    &+ 0{,}531\,\text{A}\sin(1508\,\text{rad/s}\cdot t - 0{,}879\,\text{rad})
\end{split}
\end{equation}

\textbf{Cálculo de valores RMS y componentes AC:}

\begin{equation}
    V_{rms} = \sqrt{(187\,\text{V})^2 + \frac{(125\,\text{V})^2 + (25{,}0\,\text{V})^2 + (10{,}7\,\text{V})^2}{2}} = 208\,\text{V}
\end{equation}

\begin{equation}
    I_{rms} = \sqrt{(6{,}24\,\text{A})^2 + \frac{(3{,}56\,\text{A})^2 + (0{,}531\,\text{A})^2}{2}} = 6{,}73\,\text{A}
\end{equation}

\begin{equation}
    V_{ac} = \sqrt{(208\,\text{V})^2 - (187\,\text{V})^2} = 91{,}0\,\text{V}
\end{equation}

\begin{equation}
    I_{ac} = \sqrt{(6{,}73\,\text{A})^2 - (6{,}23\,\text{A})^2} = 2{,}55\,\text{A}
\end{equation}

\textbf{Corrientes en el diodo:}

\begin{equation}
    I_{D,DC} = \frac{6{,}23\,\text{A}}{2} = 3{,}12\,\text{A}
\end{equation}

\begin{equation}
    I_{D,rms} = \frac{6{,}73\,\text{A}}{\sqrt{2}} = 4{,}76\,\text{A}
\end{equation}

\textbf{Cálculo de potencias:}

\begin{equation}
    P = (6{,}73\,\text{A})^2 \times 30{,}0\,\Omega = 1354\,\text{W}
\end{equation}

\begin{equation}
    S = 208\,\text{V} \times 6{,}73\,\text{A} = 1399\,\text{VA}
\end{equation}

\begin{equation}
    P_{DC} = 187\,\text{V} \times 6{,}23\,\text{A} = 1165\,\text{W}
\end{equation}

\textbf{Parámetros de calidad:}

\begin{equation}
    fp = \frac{1354\,\text{W}}{1399\,\text{VA}} = \boxed{0{,}968}
\end{equation}

\begin{equation}
    \eta = \frac{1165\,\text{W}}{1399\,\text{VA}} = \boxed{0{,}833}
\end{equation}

\begin{equation}
    FR_V = \frac{91{,}0\,\text{V}}{187\,\text{V}} = \boxed{0{,}487}
\end{equation}

\begin{equation}
    FR_I = \frac{2{,}55\,\text{A}}{6{,}23\,\text{A}} = \boxed{0{,}409}
\end{equation}

\subsection{Gráficas de las funciones teóricas}

A continuación se presentan las gráficas de $v_o(t)$ e $i_o(t)$ obtenidas del análisis de
Fourier (elaboradas con los datos anteriores):

\begin{figure}[H]
    \centering
    \includegraphics[width=1\textwidth]{Funvol.png}
    \caption{Función de tensión de salida $v_o(t)$ - Gráfica teórica}
\end{figure}

\begin{figure}[H]
    \centering
    \includegraphics[width=1\textwidth]{Funcor.png}
    \caption{Función de corriente de salida $i_o(t)$ - Gráfica teórica}
\end{figure}


\subsection{Gráficas de las funciones teóricas}

A continuación se presenta el código de MATLAB utilizado para generar las gráficas
de $v_o(t)$ e $i_o(t)$ a partir del análisis de Fourier obtenido en el procedimiento:

\begin{lstlisting}[caption={Código MATLAB para graficar $v_o(t)$ e $i_o(t)$ - Ejercicio 3}, label={lst:ejercicio3}]
%% Ejercicio 3 - Rectificador monofasico de onda completa tipo puente
%  Graficas de vo(t) e io(t) obtenidas del analisis de Fourier
%  a=3, b=9, c=3 -> R=30 Ohm, L=24 mH
clear; clc; close all;

%% Vector de tiempo
f0   = 60;                       % Hz
T0   = 1/f0;                     % periodo de salida
t    = linspace(0, 1*T0, 5000);  % un periodo

%% ---------- Tension de salida vo(t) ----------
VDC = 187;      % V
V1  = 125;      % V
V2  = 25.0;     % V
V3  = 10.7;     % V

w1 = 754;   % rad/s
w2 = 1508;  % rad/s
w3 = 2262;  % rad/s

vo = VDC - V1*cos(w1*t) - V2*cos(w2*t) - V3*cos(w3*t);

%% ---------- Corriente de salida io(t) ----------
IDC = 6.23;     % A
I1  = 3.57;     % A
I2  = 0.532;    % A

phi1 = 0.543;   % rad
phi2 = 0.879;   % rad

io = IDC + I1*sin(w1*t - phi1) + I2*sin(w2*t - phi2);

%% ---------- Tension de entrada (referencia) ----------
Vm = 294;       % V pico
w  = 377;       % rad/s
vin = Vm*sin(w*t);

%% ================= FIGURA 1: VOLTAJES =================
figure('Color','w','Position',[100 100 900 450]);
plot(t*1000, vin, 'r', 'LineWidth', 1.2); hold on;
plot(t*1000, vo, 'b', 'LineWidth', 1.6);
grid on; box on;
xlabel('Tiempo (ms)', 'FontSize', 11);
ylabel('Voltaje (V)', 'FontSize', 11);
title('Taller\_Punto3 - Voltajes', 'FontSize', 12, 'FontWeight', 'bold');
legend('Voltaje de entrada (v_{in})', 'Voltaje de salida (v_o)', ...
       'Location', 'southoutside', 'Orientation', 'horizontal');
set(gca, 'FontSize', 10);
xlim([0 max(t)*1000]);
exportgraphics(gcf, 'vo_ejercicio3.png', 'Resolution', 300);

%% ================= FIGURA 2: CORRIENTES =================
figure('Color','w','Position',[100 100 900 450]);
plot(t*1000, io, 'Color', [0 0.6 0], 'LineWidth', 1.6);
grid on; box on;
xlabel('Tiempo (ms)', 'FontSize', 11);
ylabel('Corriente (A)', 'FontSize', 11);
title('Taller\_Punto3 - Corriente de salida', 'FontSize', 12, 'FontWeight', 'bold');
legend('Corriente de salida (i_o)', 'Location', 'southoutside');
set(gca, 'FontSize', 10);
xlim([0 max(t)*1000]);
ylim([min(io)-0.5 max(io)+0.5]);
exportgraphics(gcf, 'io_ejercicio3.png', 'Resolution', 300);

%% ============= Valores de verificacion en consola =============
fprintf('VDC   = %.2f V\n', mean(vo));
fprintf('Vrms  = %.2f V\n', rms(vo));
fprintf('IDC   = %.2f A\n', mean(io));
fprintf('Irms  = %.2f A\n', rms(io));
\end{lstlisting}
\subsection{Simulación}

\begin{figure}[H]
    \centering
    \includegraphics[width=1\textwidth]{simulacion_ejercicio3.png}
    \caption{Simulación en Simulink del rectificador de onda completa tipo puente
    ($V_s=208\,\text{V}$, $f=60\,\text{Hz}$, $L=24{,}0\,\text{mH}$, $R=30{,}0\,\Omega$)}
\end{figure}

De la simulación se obtuvieron los siguientes valores medidos en la carga:

\begin{itemize}
    \item $I_{DC} = 6{,}07\,\text{A}$, \quad $I_{rms} = 6{,}57\,\text{A}$
    \item $V_{DC} = 182\,\text{V}$, \quad $V_{rms} = 203\,\text{V}$
\end{itemize}


\begin{figure}[H]
    \centering
    \includegraphics[width=0.9\textwidth]{simuvol.png}
    \caption{Simulación en Simulink de la salida de voltaje}
\end{figure}


\begin{figure}[H]
    \centering
    \includegraphics[width=0.9\textwidth]{simucor.png}
    \caption{Simulación en Simulink de la salida de corriente}
\end{figure}


\subsection{Análisis}

El rectificador de onda completa tipo puente con carga R-L presenta las siguientes
características fundamentales:

\textbf{Comportamiento del circuito:}
\begin{itemize}
    \item El puente rectificador convierte ambos semiciclos de la señal AC de entrada en
    pulsos positivos, duplicando la frecuencia de salida a 120\,Hz.
    \item La inductancia de 24{,}0\,mH actúa como filtro, reduciendo el rizado de corriente
    y suavizando la forma de onda.
    \item Los armónicos predominantes son pares ($m=2,4,6$) con amplitudes decrecientes
    debido al efecto de filtrado de la inductancia.
\end{itemize}

\textbf{Cálculos de valores no medidos (a partir de la simulación):}

\begin{equation}
    V_{AC} = \sqrt{(203\,\text{V})^2 - (182\,\text{V})^2} = 89{,}9\,\text{V}
\end{equation}

\begin{equation}
    I_{AC} = \sqrt{(6{,}57\,\text{A})^2 - (6{,}07\,\text{A})^2} = 2{,}51\,\text{A}
\end{equation}

\begin{equation}
    I_{D,DC} = \frac{6{,}07\,\text{A}}{2} = 3{,}04\,\text{A}
\end{equation}

\begin{equation}
    I_{D,rms} = \frac{6{,}57\,\text{A}}{\sqrt{2}} = 4{,}65\,\text{A}
\end{equation}

\begin{equation}
    P = (6{,}57\,\text{A})^2 \times 30\,\Omega = 1294\,\text{W}
\end{equation}

\begin{equation}
    S = 203\,\text{V} \times 6{,}57\,\text{A} = 1334\,\text{VA}
\end{equation}

\begin{equation}
    P_{DC} = 182\,\text{V} \times 6{,}07\,\text{A} = 1105\,\text{W}
\end{equation}

\begin{equation}
    fp = \frac{1294\,\text{W}}{1334\,\text{VA}} = 0{,}970
\end{equation}

\begin{equation}
    \eta = \frac{1105\,\text{W}}{1334\,\text{VA}} = 0{,}828
\end{equation}

\begin{equation}
    FR_V = \frac{89{,}9\,\text{V}}{182\,\text{V}} = 0{,}494
\end{equation}

\begin{equation}
    FR_I = \frac{2{,}52\,\text{A}}{6{,}07\,\text{A}} = 0{,}415
\end{equation}

\textbf{Comparación teórica vs. simulación:}
\begin{itemize}
    \item Los errores entre valores calculados y simulados se mantienen por debajo del
    5\,\% en todos los parámetros, indicando buena concordancia entre el modelo teórico
    y el circuito simulado.
    \item Las mayores diferencias se presentan en las potencias ($P$, $S$, $P_{DC}$),
    lo cual es esperable debido a la caída de tensión directa de los diodos
    (aprox.\ 0,700\,V c/u), no incluida en el modelo ideal.
    \item El factor de rizado de corriente ($\approx 0{,}409$) es menor que el de tensión
    ($\approx 0{,}487$), confirmando el efecto de filtrado de la inductancia de 24{,}0\,mH.
    \item La eficiencia de rectificación cercana al 83{,}0\,\% y el factor de potencia
    cercano a 0,968 indican un excelente desempeño del rectificador tipo puente.
\end{itemize}

\begin{table}[h]
\centering
\begin{tabular}{|l|c|c|c|}
\hline
\textbf{Parámetro} & \textbf{Valor calculado} & \textbf{Valor simulado} & \textbf{Error (\%)} \\
\hline
$t_e$ (ms)        & No aplica & No aplica & --- \\
$V_{DC}$ (V)      & 187    & 182 & 2{,}74 \\
$V_{AC}$ (V)      & 91{,}0 & 89{,}9  & 1{,}22 \\
$V_{rms}$ (V)     & 208    & 203     & 2{,}46 \\
$I_{DC}$ (mA)     & 6230   & 6070    & 2{,}64 \\
$I_{AC}$ (mA)     & 2550   & 2510    & 1{,}59 \\
$I_{rms}$ (mA)    & 6730   & 6570    & 2{,}44 \\
$I_{D,DC}$ (mA)   & 3120   & 3040    & 2{,}63\\
$I_{D,rms}$ (mA)  & 4760   & 4650    & 2{,}37 \\
$P$ (W)           & 1354   & 1294    & 4{,}64 \\
$S$ (VA)          & 1399   & 1334    & 4{,}87 \\
$fp$              & 0{,}968 & 0{,}970 & -0{,}206 \\
$P_{DC}$ (W)      & 1165   & 1105    & 5{,}43 \\
$\eta$            & 0{,}833 & 0{,}828 & 0{,}604 \\
$FR_V$            & 0{,}487 & 0{,}494 & -1{,}42 \\
$FR_I$            & 0{,}409 & 0{,}415 & -1{,}45 \\
\hline
\end{tabular}
\caption{Tabla de comparación No.3 -- Ejercicio 3 (a=3, b=9, c=3)}
\end{table}

\subsection{Resultados}

Valores principales obtenidos del rectificador de onda completa tipo puente:

\begin{enumerate}
    \item \textbf{Tensiones en la carga:}
    \begin{itemize}
        \item Tensión DC: $V_{DC} = 187\,\text{V}$
        \item Tensión RMS: $V_{rms} = 208\,\text{V}$
        \item Componente AC: $V_{ac} = 91{,}0\,\text{V}$
    \end{itemize}

    \item \textbf{Corrientes en la carga:}
    \begin{itemize}
        \item Corriente DC: $I_{DC} = 6{,}23\,\text{A} = 6230\,\text{mA}$
        \item Corriente RMS: $I_{rms} = 6{,}73\,\text{A} = 6730\,\text{mA}$
        \item Componente AC: $I_{ac} = 2{,}55\,\text{A} = 2550\,\text{mA}$
    \end{itemize}

    \item \textbf{Corrientes en el diodo:}
    \begin{itemize}
        \item Corriente DC en diodo: $I_{D,DC} = 3{,}12\,\text{A} = 3120\,\text{mA}$
        \item Corriente RMS en diodo: $I_{D,rms} = 4{,}76\,\text{A} = 4760\,\text{mA}$
    \end{itemize}

    \item \textbf{Potencias:}
    \begin{itemize}
        \item Potencia activa: $P = 1354\,\text{W}$
        \item Potencia aparente: $S = 1399\,\text{VA}$
        \item Potencia DC: $P_{DC} = 1165\,\text{W}$
    \end{itemize}

    \item \textbf{Parámetros de calidad:}
    \begin{itemize}
        \item Factor de potencia: $fp = 0{,}968$
        \item Eficiencia de rectificación: $\eta = 0{,}833 = 83{,}3\,\%$
        \item Factor de rizado de tensión: $FR_V = 0{,}487 = 48{,}7\,\%$
        \item Factor de rizado de corriente: $FR_I = 0{,}409 = 40{,}9\,\%$
    \end{itemize}
\end{enumerate}

El tiempo de extinción $t_e$ no aplica para este circuito, ya que al ser un rectificador
de onda completa tipo puente, la conducción es continua durante todo el semiciclo
positivo y negativo de la señal de entrada.

\textbf{Conclusión:} El rectificador de onda completa tipo puente con carga R-L
(R=30{,}0\,$\Omega$, L=24{,}0\,mH) presenta un excelente desempeño, con una eficiencia del
83{,}3\,\% y un factor de potencia de 0,968. La inductancia reduce efectivamente el
rizado de corriente (40,9\,\%) frente al de tensión (48,3\,\%). Los errores entre los
valores teóricos y simulados son inferiores al 5,50\,\%, validando el modelo matemático
utilizado.

\end{document}